\begin{abstract}
A new annotation scheme, language or domain usually starts with a guideline and little gold data. Prompting an LLM for the whole annotation can violate the scheme's conventions and produce malformed output, while supervised systems need thousands of examples. We observe that a guideline is mostly a procedure: most steps are deterministic, and only some call for judgment. We propose \emph{decision orchestration}, which encodes a guideline as an executable procedure separating rule nodes from model decision nodes. Code runs the steps and assembles the output; a model is consulted only at decision points, for a schema-constrained fragment or for probabilities over the options of a closed question. We then learn the procedure itself, not model weights, from a handful of gold annotations: thresholds are tuned against the task metric, uncertain decisions are routed to a stronger model, and an LLM reflector proposes edits that, for closed decisions, are kept only if they measurably improve the score. On verbal multiword expression identification (PARSEME), a closed procedure learned from 60 gold sentences, five at a time, reaches 46.3 F in English, 4.4 points above the best published result we know of and 13 above the best system of the 2018 shared task, and 57.8 F in Chinese, 18 points above an LLM annotating directly with ten gold examples; reflection costs under one dollar per language. On Uniform Meaning Representation parsing, where the output is mostly graph structure, a 9B model inside the procedure reaches 17.6 F against 4.9 for a detailed prompt, two gold documents of reflection raise a frontier model from 45.3 to 52.6 F, and routing adds 2.5 F to a closed procedure; yet direct generation with ten gold examples matches the best closed procedure in Chinese and exceeds it in English. Decision orchestration pays off most where a guideline is a sequence of tests over enumerable candidates.
\end{abstract}
